\section{Operating the Pipeline}
\label{sec:operations}

We treat operation as part of the data engineering lifecycle~\cite{reis2022fundamentals}, not as
an afterthought. Every change travels the same path. A feature branch is merged into \texttt{dev} through a pull
request reviewed by the other group member, after lint and unit tests pass on a GitHub-hosted
runner. A push to \texttt{dev} runs the pipeline against the development data area on our
self-hosted runner. A pull request from \texttt{dev} to \texttt{main} releases to production,
and the production run waits for approval from the member who did not trigger it. A weekly
scheduled run picks up FDA's monthly release; it runs code from \texttt{main} only.

Development and production differ only in configuration: which source files to ingest, where
the data lives, and Spark settings. Each run writes a quality report and a log, which GitHub
Actions shows as the run summary and keeps as artifacts.

\todo{Evidence: a scheduled production run end to end (screenshot or log excerpt), how a failure
shows up, what we monitor, and how another engineer clones and runs the project.}
